%% 58장 분할표와 비율
\chapter{Contingency Tables and Proportions}
\label{ch:contingency}

\abstract{범주 자료의 검정과 구간을 다룬다. 두 범주 변수를 교차해 센 분할표에서 주변합으로 조건을 걸면 독립 가설 아래 표의 분포가 초기하 분포가 되고 독립이 아니면 오즈비 하나에만 의존하는 비중심 초기하 분포가 된다는 것을 보인다. 이 조건부 분포에서 Fisher 정확 검정과 조건부 최대우도 오즈비, 정확한 오즈비 구간이 나온다. Pearson 카이제곱 통계량이 독립 가설 아래 자유도 $(r-1)(c-1)$의 카이제곱으로 수렴한다는 것을 사영으로 증명하고 $r\times c$ 표의 정확 검정을 계산하는 망 알고리즘을 정리한다. 조건을 걸지 않고 장애 모수에 대해 상한을 취하는 Barnard와 Boschloo의 검정은 Fisher 검정보다 검정력이 크며 Boschloo의 $p$값이 Fisher의 $p$값을 넘지 않는다는 것을 증명한다. 끝으로 두 비율의 차의 구간 넷의 정확한 포함 확률을 계산해 Wald 구간이 명목 수준 $95\%$에서 $33\%$까지 떨어지는 것을 본다.}

\section{Contingency Tables and Sampling Models}
\label{sec:sampling-models}

두 범주 변수를 교차해 센 $r\times c$ 표를 \term{분할표}{contingency table}라 한다. 칸 $(i,j)$의 개수를 $N_{ij}$, 행합을 $R_i$, 열합을 $C_j$, 총합을 $N$이라 하자. 같은 표라도 자료를 모은 방식에 따라 분포가 다르다.
\begin{itemize}
\item 다항 표집: 총합 $N$만 정해 두고 개체마다 두 변수를 잰다. 표 전체가 칸 확률 $\pi_{ij}$의 다항 분포를 따른다.
\item 곱 이항(곱 다항) 표집: 행마다 개체 수 $R_i$를 정해 두고 각 행에서 열 범주를 센다. 행마다 독립인 다항 분포이고 $2\times2$ 표에서는 두 이항 표본이다.
\item 두 주변합이 모두 정해진 표집: Fisher의 차 맛보기 실험처럼 행합과 열합을 모두 실험자가 정한다.
\end{itemize}
독립 가설 $H_0:\pi_{ij}=\pi_{i\cdot}\pi_{\cdot j}$(곱 이항이면 행마다 열 확률이 같다는 가설)는 주변 확률을 장애 모수로 남긴다. 주변합으로 조건을 걸면 그 장애 모수가 사라진다(55장의 조건부 방법).

\begin{proposition}
\label{prop:hypergeometric}
곱 이항 표집의 $2\times2$ 표에서 첫 행이 $\mathrm{Bin}(n_1,p_1)$, 둘째 행이 $\mathrm{Bin}(n_2,p_2)$를 따르고 둘이 독립이라 하자. 첫 열의 합 $S=X_1+X_2=s$로 조건을 걸면
\[
P(X_1=x\mid S=s)=\frac{\binom{n_1}x\binom{n_2}{s-x}\psi^x}{\sum_u\binom{n_1}u\binom{n_2}{s-u}\psi^u},\qquad\psi=\frac{p_1/(1-p_1)}{p_2/(1-p_2)}
\]
이다. 합은 $\max(0,s-n_2)\le u\le\min(n_1,s)$에 걸친다. 특히 $p_1=p_2$이면 $\psi=1$이고 조건부 분포는 초기하 분포 $\binom{n_1}x\binom{n_2}{s-x}/\binom{n_1+n_2}s$이다.
\end{proposition}

\begin{proof}
$P(X_1=x,X_2=s-x)=\binom{n_1}xp_1^x(1-p_1)^{n_1-x}\binom{n_2}{s-x}p_2^{s-x}(1-p_2)^{n_2-s+x}$이다. $x$에 의존하지 않는 인수 $(1-p_1)^{n_1}(1-p_2)^{n_2}\{p_2/(1-p_2)\}^s$를 빼면 남는 것은 $\binom{n_1}x\binom{n_2}{s-x}\psi^x$이다. 조건부 확률은 이것을 $x$에 대해 정규화한 것이다. $\psi=1$일 때의 정규화 상수는 Vandermonde 항등식 $\sum_u\binom{n_1}u\binom{n_2}{s-u}=\binom{n_1+n_2}s$이다.
\end{proof}

조건부 분포는 $p_1$, $p_2$를 따로 담지 않고 \term{오즈비}{odds ratio} $\psi$ 하나에만 의존한다. 이 분포를 Fisher의 \term{비중심 초기하 분포}{noncentral hypergeometric distribution}라 한다. 다항 표집에서도 행합과 열합으로 조건을 걸면 같은 분포가 나온다. $r\times c$ 표에서 독립 가설 아래 두 주변합으로 조건을 건 분포는
\begin{equation}
\label{eq:rxc-prob}
P(N_{ij}=x_{ij}\ \forall i,j)=\frac{\prod_iR_i!\prod_jC_j!}{N!\prod_{i,j}x_{ij}!}
\end{equation}
이다. 곱 다항 표집에서 행 $i$의 확률은 $R_i!\prod_j\pi_j^{x_{ij}}/x_{ij}!$이고 행들을 곱하면 $\prod_j\pi_j^{C_j}$가 표에 의존하지 않으므로 $\prod_{i,j}1/x_{ij}!$에 비례한다. 정규화 상수는 \cref{sec:network}에서 증명하는 항등식으로 $N!/(\prod_iR_i!\prod_jC_j!)$이다.

\section{Fisher's Exact Test and the Odds Ratio}
\label{sec:fisher-exact}

\subsection{The Test}

$2\times2$ 표의 독립 가설은 \cref{prop:hypergeometric}의 초기하 분포로 정확하게 검정한다. 이것이 \term{Fisher 정확 검정}{Fisher's exact test}이고 Fisher, Irwin, Yates가 1934--35년에 거의 같은 무렵 내놓았다\cite{fisher1935logic,yates1934}. 단측 $p$값은 초기하 꼬리 $P(X_1\ge x)$이다. 양측 $p$값은 55장에서 본 두 정의 가운데 관측 표보다 확률이 크지 않은 모든 표의 확률을 더하는 쪽이 표준이다. 부동소수 계산에서 확률이 같은 표를 놓치지 않도록 R은 관측 확률의 $1+10^{-7}$배 이하를 같다고 본다.

두 집단이 $10$명씩이고 첫 집단에서 $7$명, 둘째 집단에서 $2$명이 반응한 표
\[
\begin{pmatrix}7&3\\2&8\end{pmatrix}
\]
를 보자(행이 반응 여부, 열이 집단이라 하고 열합 $10$, $10$, 첫 행의 합 $9$로 조건을 건다). 왼쪽 위 칸은 $0$부터 $9$까지이고 초기하 분포가 $4$와 $5$를 가운데에 두고 대칭이다. 관측값 $7$의 확률은 $0.03215$이고 단측 $p$값은 $P(X\ge7)=0.03489$이다. 분포가 대칭이므로 양측 $p$값은 두 정의 모두 $0.06978$이다. 표본이 작고 분포가 이산이어서 이 검정의 실제 크기는 명목 수준보다 한참 작다(\cref{sec:unconditional}).

\subsection{Estimation and Intervals for the Odds Ratio}

표본 오즈비는 $\hat\psi=ad/(bc)=7\cdot8/(3\cdot2)=9.33$이다. 그 로그의 표준오차를 델타 방법(54장)으로 어림하면 $\sqrt{1/a+1/b+1/c+1/d}=1.049$이고 \term{Woolf 구간}{Woolf interval} $\exp(\log\hat\psi\pm1.96\cdot1.049)=(1.19,\,73.0)$을 얻는다\cite{woolf1955}. 이 구간은 $1$을 포함하지 않는다.

조건부 분포를 우도로 쓸 수도 있다. \cref{prop:hypergeometric}에서 $\theta=\log\psi$로 두면 조건부 확률은 $h(x)e^{\theta x}/K(\theta)$ 꼴의 지수족이고 충분 통계량이 $X_1$이다. 54장에서 본 대로 지수족의 점수는 $x-E_\theta[X_1]$이므로 \term{조건부 최대우도 추정량}{conditional maximum likelihood estimator}은 비중심 초기하 분포의 평균을 관측값에 맞추는 방정식 $E_{\hat\psi}[X_1\mid S=s]=x$의 근이다. $\log K$가 볼록이어서 평균이 $\theta$의 증가 함수이므로 관측값이 받침의 끝이 아니면 근이 하나 있다. 위의 표에서 조건부 최대우도 추정값은 $8.15$로 표본 오즈비보다 작다. 정확한 구간은 단측 정확 검정 둘을 수준 $\alpha/2$로 뒤집어 얻는다(55장). 아래 끝은 $P_\psi(X_1\ge x)=\alpha/2$, 위 끝은 $P_\psi(X_1\le x)=\alpha/2$의 근이고 둘 다 $\psi$에 대해 단조이므로 근이 하나이다. 이 표에서는 $(0.882,\,127.0)$으로 $1$을 포함한다. Fisher 검정의 양측 $p$값 $0.070$과 맞는 결론이다. 같은 자료에서 근사 구간과 정확 구간이 다른 결론을 내는 것은 표본이 작을 때 흔하다.

\section{Chi-Square Tests}
\label{sec:chi-square}

관측 개수 $O$와 귀무가설 아래의 기대 개수 $E$를 견주는 통계량
\[
X^2=\sum\frac{(O-E)^2}E,\qquad G=2\sum O\log\frac OE
\]
를 \term{Pearson 카이제곱 통계량}{Pearson's chi-squared statistic}\cite{pearson1900chi}과 우도비 통계량이라 한다. 둘은 Cressie--Read의 \term{멱 발산}{power divergence} 족
\[
2nI^\lambda=\frac2{\lambda(\lambda+1)}\sum O\Bigl[\Bigl(\frac OE\Bigr)^\lambda-1\Bigr]
\]
의 $\lambda=1$과 $\lambda\to0$의 경우이다\cite{cressie1984}. $O=E+\Delta$로 쓰고 $\sum\Delta=0$을 쓰면 $G=2\sum(E+\Delta)\log(1+\Delta/E)=\sum\Delta^2/E+O(\sum|\Delta|^3/E^2)$이다. 귀무가설 아래 $\Delta=O_p(\sqrt n)$이고 $E$는 $n$에 비례하므로 나머지는 $O_p(n^{-1/2})$이고 $G$와 $X^2$의 차는 $0$으로 확률 수렴한다.

\begin{theorem}
\label{thm:pearson-gof}
$(N_1,\dots,N_k)$가 다항 분포 $\mathrm{Mult}(n,p)$를 따르고 모든 $p_i>0$이면 $X^2=\sum_i(N_i-np_i)^2/(np_i)$는 $n\to\infty$일 때 $\chi^2_{k-1}$로 분포 수렴한다.
\end{theorem}

\begin{proof}
$Z_i=(N_i-np_i)/\sqrt{np_i}$로 두자. 다항 벡터는 독립인 범주 지시 벡터 $n$개의 합이고 그 공분산은 $n(\operatorname{diag}p-pp^{\mathsf T})$이므로 $Z$의 공분산은 $\Pi=I-uu^{\mathsf T}$, $u=(\sqrt{p_1},\dots,\sqrt{p_k})$이다. $\|u\|=1$이므로 $\Pi$는 $u$의 직교 여공간으로 가는 계수 $k-1$의 직교 사영이다. 다변수 중심극한정리(51장)로 $Z\xrightarrow{d}\Pi G$이고 $G$는 표준 정규 벡터이다. 연속 사상 정리로 $X^2=\|Z\|^2\xrightarrow{d}\|\Pi G\|^2$이고 52장의 결과로 $\|\Pi G\|^2\sim\chi^2_{k-1}$이다.
\end{proof}

독립성 검정에서는 기대 개수 $E_{ij}=R_iC_j/N$에 추정한 주변 확률이 들어 있다. 추정한 모수만큼 자유도가 준다는 것을 Fisher가 밝혔다\cite{fisher1922chi}. Pearson은 처음에 자유도를 칸의 수에서 하나 뺀 값으로 썼다.

\begin{theorem}
\label{thm:independence-df}
$r\times c$ 표가 다항 분포를 따르고 칸 확률이 $\pi_{ij}=p_iq_j$(모두 양수)이면 $X^2=\sum_{i,j}(N_{ij}-R_iC_j/N)^2/(R_iC_j/N)$은 $\chi^2_{(r-1)(c-1)}$로 분포 수렴한다.
\end{theorem}

\begin{proof}
$n=N$으로 쓰고 $W_{ij}=(N_{ij}-np_iq_j)/\sqrt{np_iq_j}$로 두면 \cref{thm:pearson-gof}의 증명처럼 $W\xrightarrow{d}(I-P\otimes Q)G$이다. 여기서 $P=aa^{\mathsf T}$, $Q=bb^{\mathsf T}$, $a=(\sqrt{p_i})$, $b=(\sqrt{q_j})$이고 칸을 $(i,j)$의 사전식 순서로 늘어놓았다. 한편
\[
\frac{R_iC_j}n-np_iq_j=q_j(R_i-np_i)+p_i(C_j-nq_j)+\frac{(R_i-np_i)(C_j-nq_j)}n
\]
이고 마지막 항은 $O_p(1)$이어서 $\sqrt n$으로 나누면 $0$으로 간다. $R_i-np_i=\sum_l\sqrt{np_iq_l}\,W_{il}$이므로
\[
\frac{q_j(R_i-np_i)}{\sqrt{np_iq_j}}=\sqrt{q_j}\sum_l\sqrt{q_l}\,W_{il}=\bigl((I\otimes Q)W\bigr)_{ij}
\]
이고 같은 방법으로 $p_i(C_j-nq_j)/\sqrt{np_iq_j}=((P\otimes I)W)_{ij}$이다. 분모의 $R_iC_j/n$을 $np_iq_j$로 바꾸는 것은 비가 $1$로 확률 수렴하므로 Slutsky 정리로 허용된다. 그러므로 $X^2$의 제곱근 벡터 $Z_{ij}=(N_{ij}-R_iC_j/n)/\sqrt{R_iC_j/n}$은
\[
Z=(I-I\otimes Q-P\otimes I)W+o_p(1)=(I-P)\otimes(I-Q)\,W-(P\otimes Q)W+o_p(1)
\]
이다. 극한에서 $(P\otimes Q)(I-P\otimes Q)G=0$이므로 둘째 항은 사라지고 $(I-P)\otimes(I-Q)$는 $P\otimes Q$와 직교하므로 $Z\xrightarrow{d}(I-P)\otimes(I-Q)G$이다. $(I-P)\otimes(I-Q)$는 계수 $(r-1)(c-1)$의 직교 사영이므로 \cref{thm:pearson-gof}의 증명처럼 $X^2=\|Z\|^2$이 $\chi^2_{(r-1)(c-1)}$로 수렴한다.
\end{proof}

$2\times2$ 표에서는 자유도가 $1$이다. Yates는 이산 분포를 연속 분포로 어림하는 오차를 줄이려고 각 칸의 $|O-E|$에서 $1/2$을 빼는 \term{연속성 보정}{continuity correction}을 제안했다\cite{yates1934}. 앞의 표에서 보정하지 않은 $X^2=5.05$의 $p$값은 $0.025$이고 보정한 $X^2=3.23$의 $p$값은 $0.072$로 Fisher 검정의 $0.070$에 가깝다. 보정한 검정은 Fisher 검정을 어림하므로 Fisher 검정만큼 보수적이다.

\section{Exact Tests for \texorpdfstring{$r\times c$}{r x c} Tables and the Network Algorithm}
\label{sec:network}

$r\times c$ 표의 Fisher 검정은 주변합이 같은 모든 표 가운데 \eqref{eq:rxc-prob}의 확률이 관측 표의 확률 이하인 표의 확률을 더한다. 행합 $(4,6,6)$, 열합 $(4,5,7)$인 표
\[
\begin{pmatrix}3&1&0\\1&3&2\\0&1&5\end{pmatrix}
\]
에서 주변합이 같은 표는 $215$개이다. 그것을 모두 세면 정확한 $p$값은 $0.0345$이다. 자유도 $4$의 카이제곱 근사는 $X^2=10.85$에서 $0.0283$, $G=12.25$에서 $0.0156$이다. 같은 $215$개의 표에서 $X^2$의 정확한 귀무 분포를 쓰면 $P(X^2\ge10.85)=0.0222$이다. 이처럼 작은 표에서는 근사들끼리도 두 배 가까이 어긋난다.

표의 수는 칸과 개수가 늘면 폭발한다. 행합 $(10,10,10)$, 열합이 모두 $6$인 $3\times5$ 표는 $196{,}651$개이다. Mehta와 Patel의 \term{망 알고리즘}{network algorithm}은 표를 하나씩 나열하지 않는다\cite{mehta1983}. 열을 왼쪽부터 하나씩 채워 간다고 하면 $k$개의 열을 채운 뒤 남은 일은 남은 행합 벡터만으로 정해진다. 이 벡터를 마디로 하는 그래프에서 표 하나는 첫 마디에서 끝 마디로 가는 경로 하나이다. 행을 서로 바꿔도 남은 일이 같으므로 남은 행합을 정렬해 마디를 합치면 마디가 더 준다. 위의 $3\times5$ 예에서 서로 다른 마디는 $221$개뿐이다. 경로를 따라 채운 칸의 $\prod1/x_{ij}!$를 곱하면 표의 확률에 비례하는 값이 된다.

\begin{lemma}
\label{lem:completion-sum}
남은 행합이 $r'_1,\dots,r'_r$, 남은 열합이 $C'_1,\dots,C'_m$이고 $n'=\sum_ir'_i=\sum_jC'_j$이면 그 주변합을 갖는 모든 표에 대한 $\prod_{i,j}1/x_{ij}!$의 합은 $n'!/(\prod_ir'_i!\prod_jC'_j!)$이다.
\end{lemma}

\begin{proof}
번호를 붙인 $n'$개의 공을 크기 $r'_1,\dots,r'_r$의 행 묶음으로 나누는 방법은 $n'!/\prod_ir'_i!$가지이다. 공들이 열에도 미리 $C'_1,\dots,C'_m$개씩 나뉘어 있다고 하면 각 나눔은 열 $j$의 공 가운데 몇 개가 행 $i$에 갔는지를 적은 표 $(x_{ij})$를 하나 정하고 표 하나에 해당하는 나눔은 $\prod_jC'_j!/\prod_ix_{ij}!$가지이다. 그러므로 $\sum_{(x_{ij})}\prod_j(C'_j!/\prod_ix_{ij}!)=n'!/\prod_ir'_i!$이고 양변을 $\prod_jC'_j!$로 나누면 된다.
\end{proof}

마디마다 남은 표에 대한 $\prod1/x!$의 최댓값(가장 긴 경로)과 최솟값(가장 짧은 경로)을 동적 계획법으로 구해 두면 가지치기를 할 수 있다(\cref{alg:network}). 지금까지 채운 부분의 값 $v$와 가장 긴 완성의 곱도 문턱 이하이면 그 아래의 모든 표가 기각역에 들므로 \cref{lem:completion-sum}의 닫힌 꼴로 한꺼번에 더한다. 가장 짧은 완성의 곱도 문턱을 넘으면 그 아래 표는 하나도 기각역에 들지 않으므로 버린다. 그 사이일 때만 다음 열로 펼친다. 같은 마디에 닿은 과거 값들 가운데 같은 값은 하나로 합쳐 나른다.

\begin{algorithm}[t]
\begin{algorithmic}[1]
\Require 행합 $R$, 열합 $C$, 관측 표의 $\prod1/x_{ij}!$에 $1+10^{-7}$을 곱한 문턱 $\tau$
\Ensure 정확한 양측 $p$값
\State 마디마다 남은 표의 $\prod1/x!$의 최댓값 $\mathrm{LP}$와 최솟값 $\mathrm{SP}$를 뒤에서부터 동적 계획법으로 구한다
\State 마디 $(0,R)$에 과거 값 $1$을 둔다; $p\gets0$; $K\gets\prod_iR_i!\prod_jC_j!/N!$
\For{$k=0,\dots,c-1$}
  \For{단계 $k$의 마디 $\nu$와 그 과거 값 $v$(겹친 수 $w$)}
    \If{$v\cdot\mathrm{LP}(\nu)\le\tau$} $p\gets p+w\,v\,K\cdot T(\nu)$ \Comment{$T$는 \cref{lem:completion-sum}의 합}
    \ElsIf{$v\cdot\mathrm{SP}(\nu)\le\tau$}
      \For{열 $k+1$의 채움 $x$(합 $C_{k+1}$, 칸마다 남은 행합 이하)}
        \State 마디 $\nu-x$에 과거 값 $v\prod_i1/x_i!$(겹친 수 $w$)를 더한다
      \EndFor
    \EndIf
  \EndFor
\EndFor
\State \Return $p$
\end{algorithmic}
\caption{$r\times c$ Fisher 정확 검정의 망 알고리즘}
\label{alg:network}
\end{algorithm}

마지막 단계의 마디는 남은 행합이 모두 $0$인 마디 하나이고 그곳의 $\mathrm{LP}=\mathrm{SP}=1$이므로 모든 경로가 둘째 줄 이전의 분기에서 처리된다. 최악의 경우 작업량은 여전히 지수적이므로 구현은 작업량에 상한을 두고 그것을 넘으면 카이제곱 근사나 몬테카를로를 권한다.

\section{Unconditional Exact Tests}
\label{sec:unconditional}

곱 이항 표집에서 실제로 정해진 것은 행합뿐인데 Fisher 검정은 열합까지 조건으로 걸어 표본 공간을 좁힌다. 그 결과 검정이 이산적이어서 보수적이다. 조건을 걸지 않고 공통 성공 확률 $\pi$를 장애 모수로 둔 채 55장의 정의대로 $p$값을 $\pi$에 대한 상한으로 잡으면
\[
p(x)=\sup_{0<\pi<1}P_\pi\bigl(T(Y)\text{가 }T(x)\text{만큼 또는 더 극단적}\bigr)
\]
이 된다. $Y$는 $\mathrm{Bin}(n_1,\pi)$와 $\mathrm{Bin}(n_2,\pi)$의 쌍이다. \term{Barnard 검정}{Barnard's test}은 $T$로 두 비율 차의 $z$ 통계량을 쓴다\cite{barnard1945}. \term{Boschloo 검정}{Boschloo's test}은 $T$로 Fisher 검정의 $p$값 자체를 쓴다\cite{boschloo1970}.

\begin{proposition}
\label{prop:boschloo-fisher}
같은 쪽의 Fisher $p$값을 $p_F$, Boschloo $p$값을 $p_B$라 하면 모든 표에서 $p_B\le p_F$이다. 그러므로 같은 수준에서 Boschloo 검정은 Fisher 검정이 기각하는 표를 모두 기각하고 검정력이 Fisher 검정보다 작지 않다.
\end{proposition}

\begin{proof}
$p_B(x)=\sup_\pi P_\pi(p_F(Y)\le p_F(x))$이다. 열합 $S$로 조건을 걸면 $p_F$는 조건부 분포에서 계산한 $p$값이므로 55장의 \cref{prop:pvalue-valid}로 모든 $s$와 $t$에서 $P(p_F(Y)\le t\mid S=s)\le t$이다. $\pi$가 무엇이든 $S$에 대해 평균하면 $P_\pi(p_F(Y)\le t)\le t$이고 $t=p_F(x)$로 두면 $p_B(x)\le p_F(x)$이다.
\end{proof}

Barnard 검정에는 이런 보장이 없지만 실제로는 대개 Fisher 검정보다 강력하다. 앞의 $2\times2$ 표에서 Barnard와 Boschloo의 $p$값은 모두 $0.0414$로 Fisher의 $0.0698$보다 작다. \cref{tab:unconditional}은 몇 가지 표본 크기에서 명목 수준 $5\%$의 실제 크기 $\sup_\pi P_\pi(\text{기각})$과 $p_1=0.8$, $p_2=0.3$에서의 검정력을 견준다. Fisher 검정의 크기는 $6$명씩일 때 $0.6\%$에 그쳐 검정력을 크게 잃는다. 보정하지 않은 카이제곱 검정은 검정력이 가장 크지만 크기가 $6.8\%$까지 올라 수준을 지키지 못한다. 두 조건 없는 검정은 수준을 지키면서 검정력을 대부분 되찾는다. $(5,15)$에서 Barnard 검정이 Boschloo 검정보다 약한 것은 \cref{prop:boschloo-fisher} 같은 보장이 Barnard 검정에 없기 때문이다.

\begin{table}[t]
\begin{sdtable}{colspec={cccccccc}, row{1,2}={font=\small\bfseries}}
\toprule
\SetCell[r=2]{c}{$(n_1,n_2)$} & \SetCell[c=4]{c}{실제 크기} & & & & \SetCell[c=3]{c}{검정력} & & \\
 & $X^2$ & Fisher & Barnard & Boschloo & Fisher & Barnard & Boschloo \\
\midrule
$(6,6)$ & $0.068$ & $0.006$ & $0.039$ & $0.039$ & $0.156$ & $0.389$ & $0.389$ \\
$(10,10)$ & $0.042$ & $0.013$ & $0.042$ & $0.042$ & $0.426$ & $0.622$ & $0.622$ \\
$(5,15)$ & $0.052$ & $0.025$ & $0.025$ & $0.039$ & $0.414$ & $0.414$ & $0.459$ \\
$(20,20)$ & $0.053$ & $0.025$ & $0.050$ & $0.046$ & $0.846$ & $0.901$ & $0.901$ \\
\bottomrule
\end{sdtable}
\caption{두 이항 표본의 양측 검정(명목 수준 $0.05$). 실제 크기는 공통 확률 $\pi$의 격자 $2$만 점에서 기각 확률의 최댓값이고 검정력은 $p_1=0.8$, $p_2=0.3$에서의 기각 확률이다. 모두 표본 공간 전체를 더해 정확히 계산했다. $X^2$는 보정하지 않은 Pearson 검정이다.}
\label{tab:unconditional}
\end{table}

\section{Intervals for the Difference of Two Proportions}
\label{sec:difference-intervals}

두 이항 비율의 차 $\delta=p_1-p_2$의 구간은 55장의 이항 구간과 같은 문제를 더 크게 겪는다. 대표적인 방법은 넷이다.
\begin{itemize}
\item Wald 구간: $\hat p_1-\hat p_2\pm z\sqrt{\hat p_1(1-\hat p_1)/n_1+\hat p_2(1-\hat p_2)/n_2}$.
\item \term{Agresti--Caffo 구간}{Agresti--Caffo interval}: 각 표본에 성공과 실패를 하나씩 더한 $\tilde p_i=(x_i+1)/(n_i+2)$로 Wald 구간을 만든다\cite{agresti2000caffo}.
\item \term{Newcombe 구간}{Newcombe interval}: 각 비율의 Wilson 구간 $(l_i,u_i)$를 엮어 $\hat\delta-\sqrt{(\hat p_1-l_1)^2+(u_2-\hat p_2)^2}$와 $\hat\delta+\sqrt{(u_1-\hat p_1)^2+(\hat p_2-l_2)^2}$을 끝으로 한다\cite{newcombe1998}.
\item \term{Miettinen--Nurminen 구간}{Miettinen--Nurminen interval}: 점수 검정 $(\hat p_1-\hat p_2-\delta)/\sqrt{\{\tilde p_1(1-\tilde p_1)/n_1+\tilde p_2(1-\tilde p_2)/n_2\}N/(N-1)}$을 뒤집는다\cite{miettinen1985}. $\tilde p_1=\tilde p_2+\delta$와 $\tilde p_2$는 $p_1-p_2=\delta$라는 제약 아래의 최대우도 추정값이다.
\end{itemize}
마지막 방법의 제약된 최대우도는 $\tilde p_2$에 대한 한 변수 문제이다. 가능한 구간 $(\max(0,-\delta),\min(1,1-\delta))$에서 로그 우도가 오목하므로 정류점이 있으면 그것이 최대점이고 정류점은 삼차 방정식의 근으로 닫힌 꼴로 쓸 수 있다. 정류점이 없으면 기울기가 가리키는 쪽 끝이 최대점이다.

$n_1=n_2=20$에서 $p_1$, $p_2$를 $0.01$부터 $0.99$까지 $99\times99$ 격자로 움직이며 정확한 포함 확률을 계산하면 \cref{tab:difference-coverage}이 된다. Wald 구간은 격자의 $41\%$에서 포함 확률이 $0.93$ 아래이고 가장 나쁜 곳에서 $0.33$이다. 나머지 셋은 평균이 명목 수준 근처이다. Newcombe 구간은 대부분에서 좋지만 비율이 $0$이나 $1$에 가까운 곳에서 $0.88$까지 내려가고 Agresti--Caffo와 Miettinen--Nurminen 구간은 최솟값도 $0.92$ 위에 머문다. $x_1=14$, $x_2=6$이면 네 구간은 차례로 $(0.116,0.684)$, $(0.088,0.639)$, $(0.090,0.619)$, $(0.087,0.641)$이다. Wald 구간만 눈에 띄게 오른쪽으로 치우친다.

\begin{table}[t]
\begin{sdtable}{colspec={lccc}, row{1}={font=\small\bfseries}}
\toprule
구간 & 최소 & 평균 & $0.93$ 아래인 격자점 \\
\midrule
Wald & $0.331$ & $0.925$ & $41.4\%$ \\
Agresti--Caffo & $0.924$ & $0.956$ & $0.1\%$ \\
Newcombe & $0.880$ & $0.952$ & $2.5\%$ \\
Miettinen--Nurminen & $0.922$ & $0.952$ & $0.1\%$ \\
\bottomrule
\end{sdtable}
\caption{$n_1=n_2=20$, 명목 수준 $95\%$인 비율 차 구간의 정확한 포함 확률. $(p_1,p_2)$를 $0.01$부터 $0.99$까지의 $99\times99$ 격자에서 계산했다.}
\label{tab:difference-coverage}
\end{table}

\begin{instardust}[\code{docs/design/htest.md} H77, H78, H79, H80, H98, H150]
\code{sd.htest.fisher\_exact}는 $r\times c$ 표에서도 언제나 정확한 망 알고리즘을 쓰고 R \code{fisher.test}의 규칙(관측 확률의 $1+10^{-7}$배 이하)을 따른다. scipy 1.18.1의 기본은 시드 없는 몬테카를로이므로 일부러 다르게 한 곳이다. 마디는 남은 행합을 정렬한 다중집합이고 가장 긴·짧은 완성은 메모 동적 계획법으로 정확히 구한다. R의 FEXACT가 로그로 가까운 과거 값을 합쳐 작은 상대오차를 들이는 것과 달리 같은 유리수만 합치고 대신 가운데 열에서 앞뒤의 목록을 만나게 한다. 작업량이 $2\times10^8$을 넘으면 \code{chi2\_contingency}를 권하는 \code{ValueError}를 낸다. \code{barnard\_exact}와 \code{boschloo\_exact}의 상한은 Taylor 꼴의 구간별 상한으로 구간을 나누며 참 상한을 상대 $2^{-48}$ 안에서 찾으므로 결과가 scipy의 격자 크기 \code{n}과 무관하다. \code{odds\_ratio}의 조건부 최대우도는 비중심 초기하 분포의 평균 방정식을 괄호를 지키는 Newton 방법으로 푼다. 두 비율 차의 점수 검정은 제약된 참 최우추정을 쓴다. statsmodels 0.15.0은 삼차식의 계수 하나가 달라 $\delta\ne0$에서 다른 근을 낸다. 개수가 $(15,40)$, $(25,45)$이고 $\delta=0.1$이면 참 추정값은 $\tilde p_2=0.42495$이다. 정확 검정은 정수 개수만 받는다.
\end{instardust}

\section{Concluding Remarks}

분할표의 정확 추론은 주변합으로 조건을 걸어 장애 모수를 없애는 데서 시작한다. 그 조건부 분포는 독립이면 초기하 분포이고 아니면 오즈비만의 비중심 초기하 분포여서 Fisher 정확 검정과 조건부 최대우도 오즈비, 정확한 오즈비 구간이 한 틀에서 나온다. 카이제곱 통계량의 자유도는 추정한 주변 확률이 사영의 차원을 줄인 결과이다. $r\times c$ 표의 정확 검정은 남은 행합을 마디로 하는 그래프에서 경로의 합으로 계산하고 가장 긴·짧은 경로로 가지를 친다. 조건을 거는 것이 보수성을 낳으므로 곱 이항 표집에서는 장애 모수에 대한 상한을 취하는 Barnard와 Boschloo의 검정이 더 강력하다. 비율 차의 구간은 Wald 구간을 피하고 점수 구간이나 더한 자료의 구간을 쓴다. 다음 장에서는 여러 검정을 한꺼번에 할 때의 오류율을 다룬다.

\chapterreferences
